\section{PRÁCTICA No. 1}

\subsection*{ABSORCIÓN DE LAS RADIACIONES GAMMA}

\textbf{OBJETIVO:} El alumno debe reconocer la diferente capacidad que tienen distintos materiales para absorber las radiaciones gamma y al mismo tiempo, darse cuenta de qué materiales son mejores que otros como blindaje contra esas mismas radiaciones.

\subsection*{MATERIAL Y EQUIPO}
\begin{itemize}
    \item 1 Detector de Radiaciones tipo Geiger-Müller.
    \item 1 Contenedor de Gammagrafía 
    \begin{itemize}
        \item[] (Conteniendo una fuente radiactiva con la suficiente actividad como para poder tener \textit{en sus costados niveles de rapidez de exposición de alrededor de 50 a 100 mR/h}.)
    \end{itemize}
    \item 3 Láminas de Madera delgada (triply) de 30 X 30 cm.
    \item 3 Láminas de Aluminio delgada de 30 X 30 cm.
    \item 3 Láminas de Fierro delgada de 30 X 30 cm.
\end{itemize}

\textbf{DESARROLLO:} Colocar el contenedor en el suelo de tal manera que uno de sus costados quede viendo hacia arriba, tomar las lecturas de rapidez de exposición a contacto del costado del contenedor y anotarlas en una libreta (mínimo tres lecturas); colocar una de las láminas de madera unida a la pared tomando nuevamente tres lecturas de contacto y anotarlas. Repetir esta secuencia con las tres láminas de madera y posteriormente, con las de aluminio y con las de fierro, por separado. Una vez teniendo anotados los datos de las lecturas obtenidas, sacar los promedios de cada una de éstas, es decir, el promedio de cuándo se usó una, dos y tres láminas de madera, de aluminio y de fierro respectivamente. Comparar los datos de cuando se colocó únicamente una lámina de madera, cuando se colocó una de aluminio y una de fierro; hacer los comentarios que crean pertinentes a este respecto y definir cuál creen que sea el mejor blindaje contra las radiaciones gamma.